\section{模型及其判断（RQ3）}
\label{sec:judgment}

模型决定了上限。从 GPT-6.1 Sol 换到 GPT-6.2 Sol，对哪些题能通过的影响比任何框架手段都大；我们写的任何文字都没有改变
模型在决策点上的选择，只有摆在它面前的事实做到了。

\para{多数失败是自信的交付}在可读题上，基线的 227 次失败中有 213 次是模型认为已经完成的交付。每一次都建立在模型
最初几分钟就定下的一种需求解读上，之后它用自己的解读来检验自己，所以它自己的检查都能通过。把同样的题与 Codex
公开的轨迹并排对比：Codex 经过一次长时间思考后定下更强的设计，再去加固它；而我们的运行在自己的检查通过后就停了。
可读题上与 Codex 的全部差距，集中在 Codex 稳定通过的 8 道题上（期望通过数 6.75 对 2.1），没有共同的机制。

\subsection{什么改变了选择}

我们截取真实试验中某个决策点的请求，改动其中一段文字后多次重放（关键节点重放）。重放与原始提示词完全一致，98\%
命中缓存，每个样本约需 13 秒。有三类输入改变了选择：
\begin{itemize}[leftmargin=1.4em]
  \item \textbf{在恰当时刻、以用户轮次给出的事实或指令。}在题目里加一行，指明某个缺陷已经确定，使下一步在 24 次重放
    中有 23 次写出了无偏估计量。第~\ref{sec:structure}~节的 GO 回答就属于这一类。
  \item \textbf{拿掉误导性的线索。}在一个人为设置的选择探针中，告诉模型代码文档里写的是哪个选项时，它在 16 次重放中
    16 次都保留了那个选项；不告诉它时，16 次都选对了。
  \item \textbf{把数据的测量结果摆到它面前。}在一道题上，这让两个 GPT 模型都从 8 次中 0 次变成 8 次中 8 次。
\end{itemize}

\subsection{什么没有改变选择}

\begin{itemize}[leftmargin=1.4em]
  \item \textbf{指令或人设中的建议}，试了五种措辞。“注释不是证据”这句话在 24 次决策中一次也没有改变；为 GPT-6.2 写的
    一段人设也没有改变任何可测的行为。
  \item \textbf{同一模型的预备步骤}，在开工前列出并决定待定的点。只对一道题起作用。
  \item \textbf{在领队旁边放另一个模型做通用阅读者或选择者。}Claude Opus 做阅读者 8 次中 0 次有效，而且起了反作用；
    在两份构建之间做选择，正确率为 0.55。另一个模型只在它知道 GPT 不知道的东西时才有帮助：在 V1 中，一个 Opus 定夺者
    把一道统计题修好了 4 次中 4 次，基线为 4 次中 0 次。
  \item \textbf{实时送达的提示。}一个适配器钩子把一段通用提示放在运行中第一次写文件的结果之后，并在之后每次请求的
    同一位置出现，因此保住了提示词缓存。作为用户轮次时，这段提示让下一步在两个截取的决策点上分别有 8 次中 8 次、8 次中
    5 到 6 次去检验关键点；作为工具结果的后缀时，分别是 8 次中 5 到 6 次和 8 次中 0 次；不加提示时 4 次中 0 次。然而在
    完整运行中，可读的 33 道题得分 21.0 对 21.0，成本相同：在它针对的 7 道开放解读型题上 $+1.0$，在其余 26 道上 $-1.0$。
    在一道题上，模型做了被要求的检查，打印出评分器期望的那种形式的分数（3.239，它自己的形式是 3.240），然后仍保留了
    自己的形式。
\end{itemize}

在截取的决策点上重放，会高估完整运行的效果：一段“先探查”的文字让一个决策从 24 次中 0 次变成 24 次中 24 次，但对完整
运行的影响是 $-0.04$（$\pm 0.85$）。恰当时刻的提示改变的是过程，不是判断；判断跟随模型自己的先验。我们因此不再把
建议当作在这些题上影响 GPT 判断的手段。

\subsection{模型版本与推理强度}

\begin{table}[t]
\centering
\small
\caption{团队 7.0、max 强度下，GPT-6.1 Sol 与 GPT-6.2 Sol 在 Terminal-Bench~4.0 上的比较。$S$ 为各题平均奖励之和
（式~\eqref{eq:score}）；“$N$ 次取最好”假设有一个完美的选择者从 $N$ 次运行中挑选。}
\label{tab:models}
\begin{tabular}{@{}lll@{}}
\toprule
指标 & GPT-6.1 Sol & GPT-6.2 Sol \\
\midrule
可读 33 道，$S$ & 约 18.8 & 约 21 \\
未读 30 道，同一批次，$S$ & 20.5 & 18.0（差 $-2.5$，$p = 0.15$） \\
未读 30 道，单次运行的平均 & 19.25（8 次运行） & 17.75（4 次运行） \\
未读 30 道，2 次取最好 / 4 次取最好 & 21.04 / 21.77 & 19.33 / 20.0 \\
多次运行结果不一致的未读题 & 8 次运行中 9/30 & 4 次运行中 6/30 \\
每次试验的模型调用次数 & 46.9 & 31.4 \\
每次试验的智能体用时（分钟，未读） & 33.7 & 28.8 \\
每次试验按公开标价的成本（未读） & 1.16 美元 & 1.15 美元 \\
\bottomrule
\end{tabular}
\end{table}

\para{从 GPT-6.1 Sol 到 GPT-6.2 Sol}表~\ref{tab:models} 在同一配置下比较了两个模型。GPT-6.2 的模型调用少三分之一，
用时少 15\%，按公开标价的单次成本相同。它能通过 GPT-6.1 很少通过的开放解读型题（某道题 5 次中 5 次，GPT-6.1 为 11 次中
0 次），却从不通过一些 GPT-6.1 有时能通过的题（8 次中 0 次，GPT-6.1 为 50 次中 17 次和 50 次中 16 次）。它也更确定：它的
瓶颈是每次运行都以同样方式做出的单个方法选择，所以多跑几次对它的帮助更小。我们在最终配置中保留了 GPT-6.2，因为未读题
上的差异在噪声范围内，而我们比较的对象用的也是 GPT-6.2。

\para{推理强度}在 V1 中，中等强度在十道缺陷题上与高强度相当（28/30 对 29/30），成本低 38\%，但在多缺陷题上落后
（0/6 对 3/6，$p \approx 0.06$）。在一组可读的 Terminal-Bench 筛查题上，之前的领队在 max 强度下 38 次试验通过 19 次，
中等强度 21 次通过 6 次，高强度 15 次通过 3 次，所以“高”并不是便宜的折中。对单领队而言，前 12 次调用用 max、之后用
中等强度，通过 15 次中 5 次，全程 max 为 15 次中 11 次，成本相同。max 强度因此成为编程任务上线时的深度，第~\ref{sec:structure}~节
的配置比较就是它的留出检验。

\para{RQ3 的回答}框架能改变模型下一步做什么（做一次检查、继续做），但改变不了它在选择点上的结论；只有摆到它面前的
事实和测量结果改变了选择。剩下的失败都是自信的单一解读，所以能动到它们的手段必须提供证据，而不是建议。
